\section{Conclusion}
\label{sec:conc}

Unlike most previous work on LLM knowledge conflicts within RAG frameworks, which focus on ``\emph{context-memory conflicts}'', our focus is on ``\emph{real-world inter-context conflicts}''. Within this setting we introduced the \texttt{WikiContradict} benchmark, which consists of 253 human-annotated instances covering different types of contradictions identified by Wikipedia editors. Our annotation of these instances, which includes isolation of relevant conflicting passages, resolution of anaphora and the creation of at least one question that highlights the contradictions, results in a benchmark that can effectively evaluate the capacity of LLMs to manage and reason over knowledge conflicts. This capacity was highlighted by the results of the evaluation of LLMs on the benchmark, in which for each contradictory instance, the LLMs were given different prompts, each one containing a different instruction and/or different conflicting information contained within the instance. Our experiments show that LLMs often struggle to correctly identify and manage real-world inter-context conflicts, indicating the usefulness of the benchmark and the need for further research in this direction. 
